\documentclass[12pt]{article}

\usepackage{fontspec}        % ضروري مع XeLaTeX
\usepackage{polyglossia}

\usepackage{color}

\usepackage{soul}  %(c) for the way one 
\usepackage{xcolor}  %(c) for the way two and three 
\usepackage[most]{tcolorbox}

% \sethlcolor{yellow}  %(c) way one
\newcommand{\Hl}[1]{\colorbox{yellow!40}{\strut #1}}  %(c) way two

% \newtcbox{\hlA}{on line,
%   colback=yellow!35,
%   colframe=yellow!35,
%   boxrule=0pt,
%   arc=2pt,
%   left=2pt, right=2pt, top=2pt, bottom=2pt,
% }


\setmainlanguage{arabic}

\newfontfamily\arabicfont[Script=Arabic,Scale=1.0]{Amiri}

\begin{document}
قال رسول الله ﷺ :

"إن أمتَكم هذه جُعِلَتْ عافيتُها في أولِها ، وإن آخرَها سيُصِيبُهم بلاءٌ وأمورٌ يُنْكِرُونها ، تَجِيءُ فتنٌ فيُدَقِّقُ بعضُها لبعضٍ ، فتَجِيءُ الفتنةُ فيقولُ المؤمنُ : هذه مُهْلِكَتي . ثم تَنْكَشِفُ ، ثم تجيءُ فيقولُ : هذه مُهْلِكَتي . ثم تَنْكَشِفُ ، فمَن أحبَّ منكم أن يُزَحْزَحَ عن النارِ ، ويُدْخَلَ الجنةَ فلتُدْرِكْه موتَتُه ، وهو مؤمنٌ باللهِ واليومِ الآخرِ ، وليَأْتِ إلى الناسِ ما يُحِبُّ أن يُؤْتَى إليه،
\Hl{\strut 
ومَن بايع إمامًا فأعطاه صفقةَ يدِه، وثمرةَ قلبِه ، فلْيُطِعْه
}
\Hl{
ما استطاع ، فإن جاء أحدٌ يُنازِعُه فاضربوا رقبةَ الآخرِ.\strut
}"

أخرجه النسائي بلفظه، وجاء في صحيح مسلم
\((1844)\)
بألفاظ مقاربة.

    % \textarabic{
    % إن أمتَكم هذه جُعِلَتْ عافيتُها في أولِها ، وإن آخرَها سيُصِيبُهم بلاءٌ وأمورٌ يُنْكِرُونها ، تَجِيءُ فتنٌ فيُدَقِّقُ بعضُها لبعضٍ ، فتَجِيءُ الفتنةُ فيقولُ المؤمنُ : هذه مُهْلِكَتي . ثم تَنْكَشِفُ ، ثم تجيءُ فيقولُ : هذه مُهْلِكَتي . ثم تَنْكَشِفُ ، فمَن أحبَّ منكم أن يُزَحْزَحَ عن النارِ ، ويُدْخَلَ الجنةَ فلتُدْرِكْه موتَتُه ، وهو مؤمنٌ باللهِ واليومِ الآخرِ ، وليَأْتِ إلى الناسِ ما يُحِبُّ أن يُؤْتَى إليه،\strut
    % }\strut \colorbox{yellow}{\textarabic{\strut
    % ومَن بايع إمامًا فأعطاه صفقةَ يدِه
    % }}
    % \colorbox{yellow}{
    %  ، وثمرةَ قلبِه ، فلْيُطِعْه ما استطاع ، فإن جاء أحدٌ يُنازِعُه فاضربوا رقبةَ الآخرِ.
    % }
    
\end{document}